% Part: first-order-logic
% Chapter: beyond
% Section: higher-order-logic

\documentclass[../../../include/open-logic-section]{subfiles}

\begin{document}

\olfileid{fol}{byd}{hol}

\olsection{उच्च-क्रम तर्कशास्त्र}

प्रथम-क्रम तर्कशास्त्रातून द्वितीय-क्रम तर्कशास्त्राकडे गेल्यामुळे औपचारिक
भाषेच्या चौकटीत प्रथम-क्रम !!{domain} मधील वस्तूंच्या संचांविषयी बोलणे शक्य
झाले. तेथेच का थांबावे? उदाहरणार्थ, तृतीय-क्रम तर्कशास्त्रामुळे वस्तूंच्या
संचांचे संच, किंवा वस्तू आणि वस्तूंचे संच हे दोन्ही अंतर्भूत करणारे संचही
हाताळता यावेत. चतुर्थ-क्रम तर्कशास्त्रामुळे अशा प्रकारच्या वस्तूंच्या
संचांविषयी बोलता येईल. तुमच्या लक्षात आले असेलच की ही कल्पना कितीही वेळा
पुनरावृत्त करता येते.

व्यवहारात उच्च-क्रम तर्कशास्त्र संबंधांऐवजी फलनांच्या मदतीने
!!{formula}ted केले जाते. (स्वाभाविक ओळखी गृहीत धरल्यास हा फरक महत्त्वाचा
नाही.) काही मूलभूत “वर्ग” $A$, $B$, $C$,~\dots (ज्यांना आता आपण “प्रकार”
म्हणू) दिले असता, पुढील अट घालून नवे प्रकार तयार करता येतात:
\begin{quote}
$\sigma$ आणि $\tau$ हे सांत प्रकार असतील, तर $\sigma \to \tau$ हाही सांत
प्रकार असतो.
\end{quote}
प्रकारांना विन्यासात्मक “लेबले” समजा; आपल्या !!{domain} मध्ये हव्या असलेल्या
वस्तूंचे ती वर्गीकरण करतात. $\sigma \to \tau$ हा प्रकार, प्रकार~$\sigma$ च्या
वस्तू घेऊन प्रकार~$\tau$ च्या वस्तू देणारी फलने असलेल्या वस्तूंचे वर्णन करतो.
उदाहरणार्थ, “सत्य” आणि “असत्य” या सत्यतामूल्यांचा $\Omega$ हा प्रकार आणि
नैसर्गिक संख्यांचा $\Nat$ हा प्रकार घ्यावासा वाटू शकतो. अशा वेळी $\Nat \to
\Omega$ प्रकारच्या वस्तू एकस्थानी संबंध किंवा $\Nat$ चे उपसंच मानता येतात;
$\Nat \to \Nat$ प्रकारच्या वस्तू ही नैसर्गिक संख्यांपासून नैसर्गिक संख्यांकडे
जाणारी फलने असतात; आणि $(\Nat \to \Nat) \to \Nat$ प्रकारच्या वस्तू
“फलनके” असतात, म्हणजे फलने घेऊन संख्या देणारी उच्च-प्रकार फलने.

द्वितीय-क्रम तर्कशास्त्राप्रमाणेच, उच्च-क्रम तर्कशास्त्राकडे बहुवर्गीय
तर्कशास्त्राचा एक प्रकार म्हणून पाहता येते; त्यात विचारात घ्यायच्या प्रत्येक
वस्तुप्रकारासाठी एक वर्ग असतो. पण उच्च-प्रकार तर्कशास्त्राचा विन्यास
मुळापासूनच परिभाषित करणे सहसा अधिक स्पष्ट ठरते. उदाहरणार्थ, सांत प्रकारांचा
संच पुढीलप्रमाणे विगमनाने परिभाषित करता येतो:
\begin{enumerate}
\item $\Nat$ हा सांत प्रकार आहे.
\item $\sigma$ आणि $\tau$ हे सांत प्रकार असतील, तर $\sigma
  \to \tau$ हाही सांत प्रकार आहे.
\item $\sigma$ आणि $\tau$ हे सांत प्रकार असतील, तर $\sigma \times
  \tau$ हाही सांत प्रकार आहे.
\end{enumerate}
सहज अर्थाने, $\Nat$ नैसर्गिक संख्यांचा प्रकार दर्शवतो, $\sigma \to \tau$
हा $\sigma$ पासून $\tau$ कडे जाणाऱ्या फलनांचा प्रकार दर्शवतो आणि
$\sigma \times \tau$ हा वस्तूंच्या जोड्यांचा प्रकार दर्शवतो, ज्यातील एक वस्तू
$\sigma$ मधील आणि दुसरी $\tau$ मधील असते. मग पदांचा संच पुढीलप्रमाणे
विगमनाने परिभाषित करता येतो:
\begin{enumerate}
\item प्रत्येक प्रकार $\sigma$ साठी $x$, $y$, $z$, \dots ही प्रकार~$\sigma$
  ची अनेक चरे असतात.
\item $\Obj 0$ हे $\Nat$ प्रकारचे पद आहे.
\item $\Obj S$ (उत्तरवर्ती) हे $\Nat \to \Nat$ प्रकारचे पद आहे.
\item $s$ हे $\sigma$ प्रकारचे पद आणि $t$ हे
  $\Nat \to (\sigma \to \sigma)$ प्रकारचे पद असेल, तर $\Obj{R}_{st}$ हे
  $\Nat \to \sigma$ प्रकारचे पद आहे.
\item $s$ हे $\tau \to \sigma$ प्रकारचे पद आणि $t$ हे प्रकार~$\tau$ चे पद
  असेल, तर $s(t)$ हे $\sigma$ प्रकारचे पद आहे.
\item $s$ हे प्रकार~$\sigma$ चे पद आणि $x$ हे प्रकार~$\tau$ चे चर असेल,
  तर $\lambd[x][s]$ हे $\tau \to \sigma$ प्रकारचे पद आहे.
\item $s$ हे प्रकार~$\sigma$ चे पद आणि $t$ हे प्रकार~$\tau$ चे पद असेल,
  तर $\tuple{s, t}$ हे $\sigma \times \tau$ प्रकारचे पद आहे.
\item $s$ हे प्रकार~$\sigma \times \tau$ चे पद असेल, तर $p_1(s)$ हे
  प्रकार~$\sigma$ चे पद आणि $p_2(s)$ हे प्रकार~$\tau$ चे पद आहे.
\end{enumerate}
सहज अर्थाने, $\Obj{R}_{st}$ हे पुढीलप्रमाणे पुनरावर्ती रीतीने परिभाषित केलेले
फलन दर्शवते:
\begin{align*}
\Obj{R}_{st}(0) & = s \\
\Obj{R}_{st}(x+1) & = t(x, R_{st}(x)),
\end{align*}
$\tuple{s, t}$ ही ज्याचा पहिला घटक~$s$ आणि दुसरा घटक~$t$ आहे अशी जोडी
दर्शवते, तर $p_1(s)$ आणि~$p_2(s)$ हे $s$ चे पहिले व दुसरे घटक
(“प्रक्षेपणे”) दर्शवतात. शेवटी, $\lambd[x][s]$ हे
\[
f(x) = s
\]
या अटीने प्रकार~$\tau$ च्या कोणत्याही~$x$ साठी परिभाषित झालेले फलन~$f$
दर्शवते.
%READERNOTE{OLFOL-025: बाब (6) मध्ये x चा प्रकार tau दिलेला असूनही गोठवलेल्या इंग्रजीतील पुढील स्पष्टीकरणात तो sigma म्हटला आहे. फलनाचा प्रांत आणि tau ते sigma हा प्रकार जुळण्यासाठी मराठीत tau केले असून गोठवलेले इंग्रजी बाइट्स बदललेले नाहीत.}
म्हणून बाब (6) आपल्याला गुणधर्माधारित फलनरचनेचे एक रूप देते आणि
पदांच्या मदतीने फलने परिभाषित करणे शक्य करते. समान प्रकारच्या पदांमधील
!!{identity} विधाने $\eq[s][t]$, नेहमीचे विधानीय संयोजक आणि उच्च-प्रकार
संख्यापन यांपासून !!^{formula}s घडवली जातात. वर परिभाषित केलेल्या पदांचे
नियमन करणारी मूलभूत समीकरणे, तसेच संख्यापक आणि !!{identity} असलेले
तर्कशास्त्राचे नेहमीचे नियम, ही मग प्रणालीची स्वयंसिद्धके मानता येतात.

सांत प्रकारांच्या प्रणालीत सत्यतामूल्यांचा $\Omega$ हा प्रकार जोडला, तर
त्याचा वापर नियंत्रित करणारी स्वयंसिद्धकेही समाविष्ट करावी लागतात. खरे तर
योग्य युक्ती वापरून गुंतागुंतीची !!{formula}s पूर्णपणे काढून टाकता येतात आणि
त्यांच्या जागी $\Omega$ प्रकारची पदे ठेवता येतात!{} त्यानुसार सिद्धता-पद्धतीतही
बदल करता येतो. त्यातून मूलतः १९३० च्या दशकात ॲलोन्झो चर्च यांनी मांडलेली
\emph{प्रकारांची साधी उपपत्ती} मिळते.

द्वितीय-क्रम तर्कशास्त्राप्रमाणेच, उच्च-प्रकार चिन्हार्थमीमांसेच्याही विविध
आवृत्त्या वापरता येतात. पूर्ण आवृत्तीत $\sigma \to \tau$ प्रकारची चरे,
प्रकार~$\sigma$ च्या वस्तूंपासून प्रकार~$\tau$ च्या वस्तूंकडे जाणाऱ्या
\emph{सर्व} फलनांच्या संचावर मान घेतात. अपेक्षेप्रमाणे ही चिन्हार्थमीमांसा
इतकी समर्थ आहे की तिच्यासाठी संपूर्ण आणि प्रभावी !!{derivation} प्रणाली
असू शकत नाही. पण दुर्बल चिन्हार्थमीमांसेचाही विचार करता येतो; त्यात
!!a{structure} मध्ये प्रत्येक प्रकार $\tau$ साठी $T_\tau$ हे घटकांचे संच
आणि उपयोजन, प्रक्षेपण इत्यादींसाठी योग्य कृत्ये असतात. तपशील योग्य रीतीने
मांडल्यास वर वर्णिलेल्या प्रकारच्या !!{derivation} प्रणालींसाठी संपूर्णता
प्रमेये मिळवता येतात.

उच्च-प्रकार तर्कशास्त्र आकर्षक आहे, कारण गणिताचा मोठा भाग स्वाभाविक रीतीने
अंतःस्थापित करता येईल अशी चौकट ते देते: $\Nat$ पासून सुरुवात करून वास्तव
संख्या, संतत फलने इत्यादी परिभाषित करता येतात. अंतःप्रज्ञावादी तर्कशास्त्राच्या
संदर्भातही ते विशेष आकर्षक आहे, कारण प्रकारांना स्पष्ट “रचनाशील” अर्थनिर्धारणे
आहेत. खरे तर, उच्च-प्रकार चिन्हार्थमीमांसेच्या रचनाशील आवृत्त्या
(अभिजात तर्कशास्त्राऐवजी अंतःप्रज्ञावादी तर्कशास्त्रावर आधारलेल्या) विकसित
करता येतात. त्या ही रचनाशील अर्थनिर्धारणे अतिशय स्पष्ट करतात आणि अनेक
दृष्टींनी त्यांच्या अभिजात समकक्षांपेक्षा अधिक रोचक असतात.

\end{document}
